% !TEX root = IE3_expG1_prelab.tex
\documentclass[lualatex,a4paper,10pt]{jlreq}
% Shared LuaLaTeX preamble.
\usepackage{luatexja-fontspec}
\setmainfont{TeX Gyre Termes}
\setsansfont{TeX Gyre Heros}
\setmonofont{Latin Modern Mono}
\setmainjfont[BoldFont=HaranoAjiGothic-Medium]{HaranoAjiMincho-Regular}
\setsansjfont[BoldFont=HaranoAjiGothic-Bold]{HaranoAjiGothic-Medium}
\usepackage[top=22mm,bottom=22mm,left=23mm,right=23mm]{geometry}
\usepackage{amsmath,amssymb,booktabs,array,longtable,tabularx}
\usepackage{graphicx,xcolor,tikz,circuitikz}
\usetikzlibrary{arrows.meta,positioning,calc}
\usepackage{enumitem,fancyhdr,listings,needspace}
\usepackage[unicode,colorlinks=true,linkcolor=labblue,urlcolor=labteal]{hyperref}
\usepackage{bookmark}
\definecolor{labblue}{RGB}{30,70,112}
\definecolor{labteal}{RGB}{0,100,105}
\definecolor{labgray}{gray}{0.95}
\ModifyHeading{section}{font={\large\sffamily\bfseries\color{labblue}}}
\ModifyHeading{subsection}{font={\normalsize\sffamily\bfseries\color{labteal}}}
\setlength{\parindent}{1\zw}
\setlength{\parskip}{3pt}
\setlength{\emergencystretch}{2em}
\renewcommand{\arraystretch}{1.35}
\setlist{itemsep=3pt,topsep=4pt,leftmargin=2\zw}
\newcolumntype{Y}{>{\raggedright\arraybackslash}X}
\newcolumntype{C}[1]{>{\centering\arraybackslash}p{#1}}
\pagestyle{fancy}
\fancyhf{}
\fancyhead[L]{\small\color{labblue} IE3 後期・電子工学実験}
\fancyhead[R]{\small テーマ\LabCode}
\fancyfoot[C]{\thepage}
\setlength{\headheight}{16pt}
\DeclareOldFontCommand{\rm}{\normalfont\rmfamily}{\mathrm}
\lstset{basicstyle=\small\ttfamily,columns=fullflexible,keepspaces=true,
 breaklines=true,frame=single,backgroundcolor=\color{labgray},
 numbers=left,numberstyle=\tiny,showstringspaces=false,aboveskip=8pt,belowskip=8pt}
\newcommand{\blank}{\rule{22mm}{0.3pt}}
\newcommand{\answerline}{\par\noindent\rule{\linewidth}{0.25pt}\par}
\newcommand{\writing}[1]{\par\Needspace{\dimexpr#1+5mm\relax}\noindent%
 \fcolorbox{labblue!35}{white}{\parbox[c][#1][t]{\dimexpr\linewidth-2\fboxsep-2\fboxrule\relax}{\vspace{2mm}\noindent\strut\hfill\par}}\par}
\newcommand{\hint}[1]{\par\smallskip\noindent{\sffamily\bfseries\color{labteal}記述の視点：}#1\par}
\newcommand{\note}[1]{\par\smallskip\noindent\fcolorbox{labblue!45}{labblue!5}{%
 \parbox{\dimexpr\linewidth-2\fboxsep-2\fboxrule\relax}{#1}}\par\smallskip}
\newcommand{\eqerror}{\varepsilon=\frac{x_{\mathrm{meas}}-x_{\mathrm{th}}}{x_{\mathrm{th}}}\times100\ \%}
\newcommand{\LabSource}{徳山工業高等専門学校 情報電子工学科，
 『2026年度 電子工学実験（3年後期）テキスト』，テーマ\LabCode，
 \expandafter\path\expandafter{\SourcePdf}。}
\newcommand{\reporttitle}{
 \noindent{\small\sffamily\color{labteal}実験報告書・記入ガイド}\par
 \noindent{\LARGE\sffamily\bfseries \LabCode\quad\LabTitle}\par\smallskip
 \noindent 氏名：\blank\quad 班：\blank\quad 実験日：\blank\par
 \note{目的・原理・手順を確認し、実測値と自分の考察を記入する。
 考察のガイドは、検討する事象・使う測定結果・記述の方針を示す。}
}
\newcommand{\prelabtitle}{
 \noindent{\small\sffamily\color{labteal}事前学習・前提知識プリント}\par
 \noindent{\LARGE\sffamily\bfseries \LabCode\quad\LabTitle}\par\smallskip
 \noindent 氏名：\blank\quad 班：\blank\quad 予習日：\blank\par
 \note{用語、回路の動き、計算、測定表への記入の順に読む。
 例題の値は説明用であり、実験結果ではない。}
}
\newcommand{\tocrow}[3]{\hyperref[#1]{#2}& #3 &\pageref*{#1}\\[2mm]}
\newcommand{\documentmap}[1]{
 \section*{目次と概要}
 \noindent 青色の項目をクリックすると該当箇所へ移動する。\par
 {\small\begin{longtable}{@{}p{0.29\linewidth}p{0.61\linewidth}r@{}}
 \toprule {\color{labblue}\bfseries 項目} & {\color{labblue}\bfseries 読むと分かること} & 頁\\\midrule
 \endhead #1\bottomrule
 \end{longtable}}
 \clearpage
}
\newenvironment{terms}{\par\smallskip\begin{longtable}{@{}p{0.23\linewidth}p{0.73\linewidth}@{}}
 \toprule {\color{labteal}\bfseries 用語}&{\color{labteal}\bfseries 意味・回路での役割}\\\midrule\endhead}
 {\bottomrule\end{longtable}}
\newcommand{\term}[2]{\textbf{#1}&#2\\[2mm]}
\newcommand{\guidepoint}[4]{
 \Needspace{32mm}\subsection{#1}
 \noindent{\bfseries\color{labteal}考える事象：}#2\par
 \noindent{\bfseries\color{labteal}参照する結果：}#3\par
 \noindent{\bfseries\color{labteal}記述の方針：}#4\par
}
\newcommand{\reportclosing}[1]{
 \clearpage\section{結論のまとめ方}\label{sec:closing}
 \hint{#1}
 \ConclusionGuide
 \begin{enumerate}
 \item 目的に対応する最も重要な実測結果を、測定条件と数値を添えて述べる。
 \item 原理と一致した点・一致しなかった点を示す。原因は確認済みの事実と推測を分ける。
 \item 実施範囲で言えることと、未確認の条件を一文ずつまとめる。
 \end{enumerate}
 \note{書き出しの型：「○○の条件で△△を測定したところ、□□となった。
 これは◇◇と整合する。ただし××は確認しておらず、一般化には追加測定が必要である。」
 空欄は自分の実測値と根拠で埋める。}
 \noindent 結論（実験後に記入）：\writing{30mm}
 \noindent 改善案・次に確認したい条件：\writing{16mm}
 \section*{参考文献}\label{sec:references}
 \begin{enumerate}
 \item \LabSource
 \item 使用したデータシート・書籍・ウェブ資料（著者、題名、該当箇所、URL、参照日）を追加する。
 \end{enumerate}
}
\newcommand{\prelabclosing}{
 \section*{準備・出典}\label{sec:prelabsource}
 \begin{itemize}[nosep]
 \item 資料、筆記具、記録用紙、電卓と、実験条件・機器設定の記録欄。
 \item 確認問題の解答と、分からない用語・回路点への具体的な質問。
 \end{itemize}
 {\small\noindent\LabSource\par}
}
\newcommand{\simpleflow}[1]{
 \begin{center}
 \begin{tikzpicture}[node distance=5mm and 5mm,
 every node/.style={font=\small},box/.style={draw=labblue,fill=labblue!4,rounded corners=1mm,
 align=center,minimum height=10mm,text width=30mm}]
 #1
 \end{tikzpicture}
 \end{center}
}

\newcommand{\LabCode}{G1}
\newcommand{\LabTitle}{マイコンによるステッピング・モータ制御}
\newcommand{\SourcePdf}{IE3_expG1_A4.pdf}
\begin{document}
\prelabtitle
\documentmap{
\tocrow{sec:auto1}{1 用語を一つずつ理解する}{言葉の意味、単位、回路や測定での役割を確認する。}
\tocrow{sec:auto2}{2 出力ビットから回転までをたどる}{部品と信号の経路を追い、状態が変化する理由を理解する。}
\tocrow{sec:auto3}{3 ステッピング・モータの前提}{この項目の仕組みと実験で確認すべき点を整理する。}
\tocrow{sec:auto4}{4 ビットと出力順序}{励磁配線と16進出力列を対応させる。}
\tocrow{sec:auto5}{5 回転数と遷移数}{一周・90度の回数と速度を計算する。}
\tocrow{sec:auto6}{6 プログラム設計の注意}{初期状態とCALL・遷移回数の違いを確認する。}
\tocrow{sec:auto7}{7 測定からレポートへ：角度と回数を合わせる}{実測値から計算・作図・記述へ進む順序を例題で練習する。}
\tocrow{sec:auto8}{8 確認問題}{自分で計算・説明して、実験に必要な理解を確かめる。}
\tocrow{sec:auto9}{解答の目安}{数値と理由を照合し、単位や論理の取り違えを見直す。}
\tocrow{sec:auto10}{準備する記録}{接続・記録欄・必要な質問を準備する。}
\tocrow{sec:prelabsource}{準備・出典}{持参物と資料を確認し、具体的な質問を準備する。}
}

\section{用語を一つずつ理解する}\label{sec:auto1}
\begin{terms}
\term{励磁・相}{コイルへ電流を流して磁界を作ること。相は制御するコイルのまとまり。出力コードの各ビットとコイルの対応を先に確かめる。}
\term{ステップ・遷移}{励磁状態を次へ移して回転子が進む操作。初期状態を設定した時点と、その後の一遷移を区別する。}
\term{1相・1--2相}{一相だけを順に励磁する方式と、一相・二相を交互に励磁する方式。後者は中間の位置を作り、角度刻みを細かくする。}
\term{ステップ角}{一遷移に対応する角度。モータと励磁方式で決まる。一周の実測遷移数から求め、原理図の模式的な角度をそのまま使わない。}
\term{脱調}{指令した遷移に回転子が追従せず、想定の角度からずれること。指令パルスが出ているだけでは正しい回転を保証しない。}
\term{保持・負荷}{停止時にも励磁して位置を保つ場合がある。負荷や慣性は加速・高速動作に影響し、脱調する間隔を変える。}
\term{ビット・16進数}{出力の各0/1がコイル状態を表す。16進数1桁は4ビット。05hは0101であり、数字5度を指示するものではない。}
\end{terms}
\section{出力ビットから回転までをたどる}\label{sec:auto2}
\subsection{配線対応が励磁の順序を決める}
資料では$D_0=A,D_2=B,D_1=\overline A,D_3=\overline B$。
したがってA、B、$\overline A$、$\overline B$の順は01、04、02、08となる。
単純に01、02、04、08とビット順に出力すれば、物理的な相の順序が変わってしまう。
反対向きへ回すには、同じ励磁位置を逆順にたどる。
\subsection{1--2相が中間位置を作る}
01hでA一相、05hでAとBの二相、04hでB一相へ進む。
二相の磁界が一相同士の間の位置を作るので、一相から次の一相までを二遷移に分ける。
同じ一周でも遷移数は倍になるが、同じパルス間隔で回転速度も同じになるわけではない。
\subsection{TeCと駆動回路の役割}
TeCのOUT命令がビット列を出力し、実験ボードの駆動回路がコイルへ必要な電流を流す。
マイコンの出力から直接大きなコイル電流を供給していると考えない。
0Chの方向設定と07hの励磁出力を区別し、指定のケーブルとボードを用いる。

\section{ステッピング・モータの前提}\label{sec:auto3}
励磁を一つ進めると一定角度だけ回る。1相励磁は一つの相を、1--2相励磁は一相と二相を交互に励磁する。資料にある原理説明用の角度を、使用モータの実測ステップ角と取り違えない。
\section{ビットと出力順序}\label{sec:auto4}
資料の配線は$D_0=A,D_2=B,D_1=\overline A,D_3=\overline B$であり、単純なビット順に相を並べると回転順序が変わる。
\par\noindent\begin{tabularx}{\linewidth}{lY}
\toprule 方式 & 16進コード\\\midrule
1相 & 01 $\to$ 04 $\to$ 02 $\to$ 08 $\to$ 01\\
1--2相 & 01 $\to$ 05 $\to$ 04 $\to$ 06 $\to$ 02 $\to$ 0A $\to$ 08 $\to$ 09 $\to$ 01\\\bottomrule
\end{tabularx}
例：05hは00000101で、AとBを同時に励磁する。逆転は状態列を逆順にたどる。
\section{回転数と遷移数}\label{sec:auto5}
\[
 \theta=\frac{360^\circ}{N},\qquad n_{90}=\frac N4,\qquad
 \text{回転速度[rpm]}=\frac{60}{N\,T_s}.
\]
仮に1相でN=200、$T_s=5\,\mathrm{ms}$なら1.8度/遷移、90度は50遷移、60 rpm。
1--2相でN=400なら0.9度、90度は100遷移となる。これらは説明例で、実際のNは確認する。
\clearpage
\section{プログラム設計の注意}\label{sec:auto6}
\simpleflow{
\node[box] (a) {初期状態を出力};
\node[box,right=of a] (b) {次状態を出力\\待ち時間};
\node[box,right=of b] (c) {遷移数を更新};
\node[box,right=of c] (d) {指定数で停止};
\draw[-{Stealth}] (a)--(b);\draw[-{Stealth}] (b)--(c);\draw[-{Stealth}] (c)--(d);
}
初期状態を設定した時点を基準位置にする。サンプルのサブルーチンが4状態をまとめて進めるなら、CALL回数とステップ遷移数を区別する。
待ち時間を急に短くすると回転子が追従できなくなることがある。
\section{測定からレポートへ：角度と回数を合わせる}\label{sec:auto7}
\subsection{一周を基準にする}
低速で同じ位置に戻るまでの\textbf{状態遷移}を数える。
初期出力を0回目とし、途中の同じ励磁コードが一周期で戻っても機械的に一周とは限らないことに注意する。
一周N遷移なら90度はN/4遷移。N/4が整数でなければ、その方式の角度刻みでは90度ちょうどを表せない。
\subsection{CALL回数とステップを分ける例}
説明用にN=48なら90度は12遷移。
一回のサブルーチンが4遷移を進める場合、3 CALLが必要である。
1--2相でN=96なら90度は24遷移。同じCALL内の遷移数かどうかをコードで再確認する。
これらのNは仮の例であり、使用モータの実測値を表へ記入する。
\subsection{速度の記録}
$\mathrm{rpm}=60/(NT_s)$で、N=200、間隔10 msの例は30 rpm。
間隔を半分にしても、脱調が起これば実際の速度や角度は計算通りにならない。
表には設定間隔、実際の回転状態、角度のずれ、負荷条件を並べる。

\clearpage
\section{確認問題}\label{sec:auto8}
\begin{enumerate}
\item 0Ahで励磁される二つの相を答える。\answerline
\item 1相で一周48遷移だった場合の90度の遷移数を求める。\answerline
\item 一回のCALLで4遷移進む場合、上の90度には何CALL必要か。\answerline
\item 1--2相の09hの次に順方向で出力するコードを答える。\answerline
\end{enumerate}
\section*{解答の目安}\label{sec:auto9}
(1) $\overline A,\overline B$。(2) 12遷移。(3) 3 CALL（初期設定と終了状態を別途確認）。
(4) 01h。
\section*{準備する記録}\label{sec:auto10}
配線対応表、励磁列、方向、90度に必要な回数、遅延ループ、実際の角度、脱調した条件をノートに用意する。
\prelabclosing
\end{document}
